% ---------------------------------------------------------------------
% Argomento: conversione binario -> decimale (somma delle potenze di 2)
%   4 numeri per fila (8, 8, 9 e 10 bit), un solo numero pari
%   fila A: 10110111=183, 101100101=357, 11010110=214, 1011000101=709
%   fila B: 110010011=403, 10100111=167, 1010100101=677, 11101100=236
%   fila C: 11001101=205, 1010001101=653, 101001011=331, 10110010=178
% Livello: base   |   Tempo stimato: 10 min
% ---------------------------------------------------------------------
\begin{esercizio}[8]
Converti in decimale:
\ABC{%
\par\smallskip\noindent\hspace*{1em}\begin{tabular}{@{}p{0.42\linewidth}p{0.42\linewidth}@{}}
  \textbf{a)}~$10110111_{2} = \ldots\ldots\ldots_{10}$ & \textbf{b)}~$101100101_{2} = \ldots\ldots\ldots_{10}$\\[4pt]
  \textbf{c)}~$11010110_{2} = \ldots\ldots\ldots_{10}$ & \textbf{d)}~$1011000101_{2} = \ldots\ldots\ldots_{10}$
\end{tabular}}{%
\par\smallskip\noindent\hspace*{1em}\begin{tabular}{@{}p{0.42\linewidth}p{0.42\linewidth}@{}}
  \textbf{a)}~$110010011_{2} = \ldots\ldots\ldots_{10}$ & \textbf{b)}~$10100111_{2} = \ldots\ldots\ldots_{10}$\\[4pt]
  \textbf{c)}~$1010100101_{2} = \ldots\ldots\ldots_{10}$ & \textbf{d)}~$11101100_{2} = \ldots\ldots\ldots_{10}$
\end{tabular}}{%
\par\smallskip\noindent\hspace*{1em}\begin{tabular}{@{}p{0.42\linewidth}p{0.42\linewidth}@{}}
  \textbf{a)}~$11001101_{2} = \ldots\ldots\ldots_{10}$ & \textbf{b)}~$1010001101_{2} = \ldots\ldots\ldots_{10}$\\[4pt]
  \textbf{c)}~$101001011_{2} = \ldots\ldots\ldots_{10}$ & \textbf{d)}~$10110010_{2} = \ldots\ldots\ldots_{10}$
\end{tabular}}

\begin{solution}
\ABC{%
  \textbf{a)}~$10110111_2 = 128 + 32 + 16 + 4 + 2 + 1 = 183_{10}$\\[4pt]
  \textbf{b)}~$101100101_2 = 256 + 64 + 32 + 4 + 1 = 357_{10}$\\[4pt]
  \textbf{c)}~$11010110_2 = 128 + 64 + 16 + 4 + 2 = 214_{10}$\\[4pt]
  \textbf{d)}~$1011000101_2 = 512 + 128 + 64 + 4 + 1 = 709_{10}$\\[4pt]
  Si sommano i pesi $2^0, 2^1, 2^2, \ldots$ (da destra) dei soli bit a 1.}{%
  \textbf{a)}~$110010011_2 = 256 + 128 + 16 + 2 + 1 = 403_{10}$\\[4pt]
  \textbf{b)}~$10100111_2 = 128 + 32 + 4 + 2 + 1 = 167_{10}$\\[4pt]
  \textbf{c)}~$1010100101_2 = 512 + 128 + 32 + 4 + 1 = 677_{10}$\\[4pt]
  \textbf{d)}~$11101100_2 = 128 + 64 + 32 + 8 + 4 = 236_{10}$\\[4pt]
  Si sommano i pesi $2^0, 2^1, 2^2, \ldots$ (da destra) dei soli bit a 1.}{%
  \textbf{a)}~$11001101_2 = 128 + 64 + 8 + 4 + 1 = 205_{10}$\\[4pt]
  \textbf{b)}~$1010001101_2 = 512 + 128 + 8 + 4 + 1 = 653_{10}$\\[4pt]
  \textbf{c)}~$101001011_2 = 256 + 64 + 8 + 2 + 1 = 331_{10}$\\[4pt]
  \textbf{d)}~$10110010_2 = 128 + 32 + 16 + 2 = 178_{10}$\\[4pt]
  Si sommano i pesi $2^0, 2^1, 2^2, \ldots$ (da destra) dei soli bit a 1.}
  \criterioABC
  {2 pt per ogni numero: 1 pt pesi (potenze di 2) corretti, 1 pt somma e risultato. Risultati: a) $183_{10}$; b) $357_{10}$; c) $214_{10}$; d) $709_{10}$.}
  {2 pt per ogni numero: 1 pt pesi (potenze di 2) corretti, 1 pt somma e risultato. Risultati: a) $403_{10}$; b) $167_{10}$; c) $677_{10}$; d) $236_{10}$.}
  {2 pt per ogni numero: 1 pt pesi (potenze di 2) corretti, 1 pt somma e risultato. Risultati: a) $205_{10}$; b) $653_{10}$; c) $331_{10}$; d) $178_{10}$.}
\end{solution}
\end{esercizio}
